% ===========================================================================
% Báo cáo Final: ViMamba-SER
% Format : Springer LNCS
% Compiler: XeLaTeX
% ===========================================================================

\documentclass[runningheads]{llncs}

\usepackage{fontspec}
\usepackage{polyglossia}
\setmainlanguage{vietnamese}
\setotherlanguage{english}

\setmainfont{Times New Roman}
\setsansfont{Arial}
\setmonofont{Courier New}

\usepackage{graphicx}
\graphicspath{{figures/}}
\usepackage{amsmath,amssymb}
\usepackage{booktabs}
\usepackage{multirow}
\usepackage{xcolor}
\usepackage{placeins}
\usepackage{array}
\usepackage{float}

\usepackage{tikz}
\usetikzlibrary{shapes.geometric, arrows.meta, positioning, fit, backgrounds}

\usepackage[colorlinks=true, citecolor=blue, linkcolor=blue,
            urlcolor=blue, bookmarks=false]{hyperref}

% ===========================================================================
\begin{document}

\title{ViMamba-SER: Nhận dạng Cảm xúc Giọng nói Tiếng Việt\\
bằng Kiến trúc Lai và Text-Enhancement}

\titlerunning{ViMamba-SER: Vietnamese SER với Text-Enhancement}

\author{Ngô Quang Huy\inst{1} \and
        Võ Tấn Đạt\inst{1} \and
        Lý Huỳnh Khánh\inst{1}}
\authorrunning{N.Q. Huy et al.}
\institute{Khoa Trí tuệ Nhân tạo, Trường Đại học FPT,\\
TP. Hồ Chí Minh, Việt Nam\\
\email{SE190631@fpt.edu.vn}}

\maketitle

% ===========================================================================
\begin{abstract}
Báo cáo này trình bày toàn bộ quá trình xây dựng hệ thống nhận dạng cảm xúc
giọng nói tiếng Việt (Vietnamese Speech Emotion Recognition -- SER), từ giai
đoạn khảo sát baseline đến mô hình cuối cùng với ứng dụng Web Demo.
Hệ thống sử dụng WavLM-Base làm bộ mã hóa âm thanh, PhoWhisper sinh
transcript tự động, PhoBERT-v2 mã hóa văn bản, module Text-aware Modality
Enhancement (TME) căn chỉnh hai luồng đặc trưng bằng cross-attention và
cosine gate, và mạng BiGRU hai chiều làm tầng fusion.
Trên tập dữ liệu ViSEC gồm 5.280 mẫu và 4 nhãn cảm xúc, nhóm thực hiện
chuỗi thí nghiệm có hệ thống: audio-only baseline đạt 56,44\%,
simple fusion cải thiện không đáng kể (tốt nhất 57,95\%), trong khi kiến trúc
TME + BiGRU kết hợp Focal Loss đạt 56,16\% nhưng cân bằng hơn giữa các lớp,
với lớp Sad tăng từ 53,02\% lên 58,14\%.
Toàn bộ hệ thống được đóng gói thành ứng dụng Web Demo (Gradio) cho phép
nhận dạng cảm xúc real-time từ microphone.

\keywords{Speech Emotion Recognition \and Vietnamese \and WavLM \and
PhoBERT \and PhoWhisper \and BiGRU \and Focal Loss \and TME \and ViSEC}
\end{abstract}

% ===========================================================================
\section{Giới thiệu}
\label{sec:intro}

Nhận dạng cảm xúc giọng nói (SER) là bài toán phân loại trạng thái cảm xúc
của người nói từ tín hiệu âm thanh, với ứng dụng rộng rãi trong chăm sóc
sức khỏe tâm thần, giao tiếp người--máy và giám sát trung tâm cuộc
gọi~\cite{vnemos2024}.
Tiếng Việt, với hệ thống sáu thanh điệu, đặt ra thách thức đặc thù:
thanh điệu thay đổi cao độ tương tự với biến đổi cảm xúc, khiến các đặc
trưng prosody truyền thống dễ bị nhầm lẫn.

Các hướng nghiên cứu SER gần đây khai thác bổ sung thông tin văn bản bên
cạnh âm thanh (multimodal fusion) và sử dụng kiến trúc sequence modeling
hiệu quả như Mamba~\cite{gu2023mamba}.
TF-Mamba~\cite{tfmamba2025} là một ví dụ kết hợp cả hai hướng này cho
sentiment analysis đa phương thức.
Tuy nhiên, tại thời điểm khảo sát, chưa tìm thấy công trình nào áp dụng
kiến trúc tương tự cho SER tiếng Việt.

Đề tài ban đầu hướng đến sử dụng Bi-Mamba làm tầng fusion. Trong quá trình
triển khai, nhóm nhận thấy hai giới hạn thực tế: (1)~quy mô dữ liệu chưa
đủ lớn cho kiến trúc phức tạp; (2)~thư viện \texttt{mamba-ssm} yêu cầu
CUDA $\geq$ 11.8, khó cài đặt trên Windows và Colab miễn phí.
Vì vậy, nhóm thiết kế cơ chế Fallback sang BiGRU, kết hợp giảm chiều ẩn
(Model Compression) và Focal Loss~\cite{lin2017focal} để tối ưu mô hình
cho dữ liệu nhỏ.

Bảng~\ref{tab:gap} tóm tắt các khoảng trống nghiên cứu được xác định.

\begin{table}[H]
\centering
\caption{Tổng hợp khoảng trống nghiên cứu.}
\label{tab:gap}
\begin{tabular}{@{}p{3.4cm}p{5.8cm}c@{}}
\toprule
\textbf{Khoảng trống} & \textbf{Mô tả} & \textbf{Mức độ} \\
\midrule
Sequence modeling nâng cao cho SER tiếng Việt
  & Chưa có công trình áp dụng Mamba hoặc BiGRU fusion cho ngôn ngữ thanh điệu & Cao \\[3pt]
Text-enhancement cho tiếng Việt
  & ASR tiếng Việt (PhoWhisper) chưa được tích hợp vào pipeline SER & Cao \\[3pt]
Alignment hai luồng đặc trưng
  & Fusion đơn giản (concatenation) không tận dụng được quan hệ thời gian giữa audio--text & Cao \\
\bottomrule
\end{tabular}
\end{table}

% ===========================================================================
\section{Dữ liệu}
\label{sec:data}

\subsection{ViSEC}

ViSEC~\cite{visec2024} là tập dữ liệu SER tiếng Việt công khai trên
HuggingFace (\texttt{hustep-lab/ViSEC}), gồm 5.280 đoạn audio 16kHz,
4 nhãn cảm xúc: \textit{happy} (1.228), \textit{neutral} (1.507),
\textit{sad} (1.079), \textit{angry} (1.466).
Phân bố nhãn tương đối cân bằng; phần lớn clip có độ dài dưới 10 giây.
Tập dữ liệu không cung cấp transcript sẵn.

\subsection{Chất lượng Transcript từ PhoWhisper}

PhoWhisper-Base~\cite{phowhisper2023} sinh transcript cho toàn bộ 5.280 mẫu.
Kết quả phân tích: 0\% transcript rỗng, 4,3\% rất ngắn ($\leq$2 từ), độ dài
trung bình 8,8 từ/đoạn.
Độ dài theo nhãn không khác biệt đáng kể (Bảng~\ref{tab:transcript}), xác
nhận PhoWhisper trích xuất nội dung lời nói nhưng không bảo tồn thông tin
prosody mang tính cảm xúc.

\begin{table}[H]
\centering
\caption{Thống kê chất lượng transcript PhoWhisper trên ViSEC.}
\label{tab:transcript}
\begin{tabular}{@{}lrrr@{}}
\toprule
\textbf{Nhãn} & \textbf{Mean (từ)} & \textbf{Min} & \textbf{Max} \\
\midrule
angry   & 8,4 & 1 & 222 \\
happy   & 9,7 & 1 & 224 \\
neutral & 9,3 & 1 & 223 \\
sad     & 7,7 & 1 & 418 \\
\midrule
\textbf{Overall} & \textbf{8,8} & 1 & 418 \\
\bottomrule
\end{tabular}
\end{table}

% ===========================================================================
\section{Phương pháp Đề xuất}
\label{sec:method}

\subsection{Tổng quan Kiến trúc}

Hệ thống ViMamba-SER gồm bốn thành phần chính (Hình~\ref{fig:pipeline}):
(1)~WavLM-Base~\cite{wavlm2022} trích xuất chuỗi đặc trưng audio $(T_a, 768)$;
(2)~PhoWhisper~\cite{phowhisper2023} sinh transcript, PhoBERT-v2~\cite{phobert2020}
mã hóa thành chuỗi $(T_t, 768)$;
(3)~module TME căn chỉnh hai luồng;
(4)~BiGRU fusion + MLP classifier.
Toàn bộ pretrained model được giữ frozen; chỉ huấn luyện TME, Projection,
BiGRU và MLP.

\begin{figure}[H]
\centering
\includegraphics[width=0.85\textwidth]{ViMamba.png}
\caption{Pipeline tổng thể ViMamba-SER. Các pretrained model (WavLM,
PhoWhisper, PhoBERT) được giữ frozen; chỉ huấn luyện phần trong khung nét đứt.}
\label{fig:pipeline}
\end{figure}

\subsection{Text-aware Modality Enhancement (TME)}
\label{sec:tme}

Module TME lấy cảm hứng từ TF-Mamba~\cite{tfmamba2025}, gồm hai bước:

\paragraph{Cross-attention alignment.}
Audio embedding đóng vai trò query, text embedding đóng vai trò key/value
trong multi-head cross-attention (8 heads, dropout 0,1).
Kết quả là chuỗi \textit{text\_aligned} có cùng chiều thời gian $T_a$ với
audio, mang thông tin văn bản đã được căn chỉnh theo từng frame audio.

\paragraph{Cosine gate.}
Độ tương đồng cosine giữa audio gốc và text\_aligned tại từng vị trí thời
gian đi qua Linear + Sigmoid tạo gate $g \in [0, 1]$:
\begin{equation}
\mathbf{e}_t = \mathbf{a}_t + g_t \cdot \mathbf{t}^{aligned}_t
\label{eq:tme}
\end{equation}
Gate cho phép mô hình tự quyết định mức độ bổ sung text: nếu text không
liên quan ($g_t \approx 0$), audio gốc được giữ nguyên.

\subsection{BiGRU Fusion và Model Compression}
\label{sec:fusion}

Chuỗi enhanced từ TME có chiều 768. Trước khi đưa vào BiGRU, một lớp
Linear Projection giảm chiều từ 768 xuống \textbf{256}. BiGRU 2 lớp (hidden
size = 128 mỗi chiều, output = 256) xử lý chuỗi hai chiều.
Sau BiGRU, concat mean-pool và max-pool tạo vector $(B, 512)$ đưa vào
MLP classifier 3 lớp (512$\to$ReLU$\to$Dropout(0,3)$\to$256$\to$ReLU$\to$4).

Giảm chiều mang lại hai lợi ích: (1)~giảm tham số cần học, chống overfit
trên tập 5.280 mẫu; (2)~giảm dung lượng checkpoint khoảng 80MB, tăng tốc
inference cho Web Demo.

\subsection{Focal Loss}
\label{sec:focal}

Để giải quyết mất cân bằng hiệu suất per-class (lớp Sad luôn thấp nhất
trong Midterm), nhóm áp dụng Focal Loss~\cite{lin2017focal}:
\begin{equation}
FL(p_t) = -\alpha_t (1 - p_t)^\gamma \log(p_t)
\end{equation}
với $\gamma = 2.0$ và $\alpha = [1{,}0;\; 1{,}5;\; 2{,}5;\; 1{,}0]$ lần
lượt cho Happy, Neutral, Sad, Angry.
Hệ số $\alpha_{Sad} = 2{,}5$ ép mô hình tập trung phạt nặng khi phân loại
sai mẫu Sad.

% ===========================================================================
\section{Kết quả Thực nghiệm}
\label{sec:results}

\subsection{Cấu hình chung}

Chia dữ liệu: 5-fold Stratified K-Fold (seed=42), trong mỗi fold tách
tiếp train/val = 85/15.
Optimizer Adam (weight\_decay=1e-4).
Phase A/B (Midterm): lr=1e-3, 50 epochs.
Phase F (Final): lr=5e-4, 80 epochs, Focal Loss.
Kết quả báo cáo: test accuracy tại epoch có val accuracy tốt nhất.

\subsection{Giai đoạn Midterm: Baseline}
\label{sec:midterm}

Giai đoạn Midterm trả lời hai câu hỏi: WavLM embedding có đủ mạnh cho SER
tiếng Việt không, và thêm text từ PhoWhisper có cải thiện không.
Hình~\ref{fig:midterm_arch} minh họa kiến trúc hai baseline.

\begin{figure}[H]
\centering
\includegraphics[width=0.90\textwidth]{mid.png}
\caption{Kiến trúc baseline Midterm.
Trái: Phase A (audio-only, WavLM mean-pool $\to$ MLP).
Phải: Phase B (audio+text fusion $\to$ MLP).}
\label{fig:midterm_arch}
\end{figure}

\paragraph{Phase A: Audio-only.}
Bảng~\ref{tab:phase_a} trình bày ablation study.
MLP 2 lớp (A1) là baseline tốt nhất với 56,44\%.
BatchNorm (A4) gây giảm do batch statistics không ổn định trên tập nhỏ.

\begin{table}[H]
\centering
\caption{Phase A: Ablation study audio-only. WavLM-Base frozen.}
\label{tab:phase_a}
\begin{tabular}{@{}llr@{}}
\toprule
\textbf{ID} & \textbf{Mô hình} & \textbf{Test Acc.} \\
\midrule
A1 & WavLM $\to$ MLP 2 lớp (baseline)           & 56,44\% \\
A2 & WavLM $\to$ Linear 1 lớp                    & 53,79\% \\
A3 & WavLM $\to$ MLP + Dropout(0.3)              & 56,06\% \\
A4 & WavLM $\to$ MLP + BatchNorm                 & 52,27\% \\
A5 & WavLM $\to$ MLP (mean $\pm$ std, 3 seeds)   & 53,45 $\pm$ 1,79\% \\
\bottomrule
\end{tabular}
\end{table}

\paragraph{Phase B: Simple Fusion.}
Bảng~\ref{tab:phase_b} so sánh các chiến lược fusion.
Text-only (B1: 44,32\%) kém audio-only 12\%, xác nhận cảm xúc tiếng Việt
nằm chủ yếu ở prosody.
Weighted fusion $\alpha{=}0{,}9$ (B3b) đạt tốt nhất 57,95\%, nhưng cải
thiện chỉ +1,51\%.

\begin{table}[H]
\centering
\caption{Phase B: So sánh chiến lược fusion. \textbf{Đậm}: tốt nhất.}
\label{tab:phase_b}
\begin{tabular}{@{}llr@{}}
\toprule
\textbf{ID} & \textbf{Chiến lược} & \textbf{Test Acc.} \\
\midrule
A1  & Audio-only (tham chiếu)                     & 56,44\% \\
B1  & Text-only (PhoBERT từ transcript)            & 44,32\% \\
B2  & Concat fusion $\to$ MLP                      & 56,57\% \\
B3  & Weighted: $\alpha{=}0{,}8$ audio             & 56,82\% \\
\textbf{B3b} & \textbf{Weighted: $\alpha{=}0{,}9$ audio} & \textbf{57,95\%} \\
B5  & Late fusion: voting                          & 57,45\% \\
\bottomrule
\end{tabular}
\end{table}

\subsection{Giai đoạn Final: TME + BiGRU + Focal Loss}
\label{sec:final}

Kết quả Midterm cho thấy simple fusion không khai thác được text hiệu quả.
Giai đoạn Final triển khai kiến trúc đầy đủ: TME + BiGRU 256 chiều + Focal
Loss (cấu hình: seed 42, lr=5e-4, 80 epochs).

\subsubsection{Kết quả Overall.}
Mô hình hội tụ tại epoch 61 (Val Acc: 54,42\%), đạt Test Accuracy
\textbf{56,16\%}.
Tổng điểm giảm nhẹ so với mô hình BiGRU 768d không Focal Loss (59,09\%),
nhưng phân bố per-class thay đổi theo đúng mục tiêu thiết kế.

\subsubsection{Per-class Trade-off.}
Bảng~\ref{tab:final_results} là kết quả quan trọng nhất của báo cáo.

\begin{table}[H]
\centering
\caption{So sánh per-class accuracy: BiGRU 768d + CrossEntropy (Trước) vs.
BiGRU 256d + Focal Loss (Sau). Mục tiêu: tăng Sad và Neutral.}
\label{tab:final_results}
\begin{tabular}{@{}lccc@{}}
\toprule
\textbf{Cảm xúc} & \textbf{Trước (768d, CE)} & \textbf{Sau (256d, Focal)} & \textbf{$\Delta$} \\
\midrule
Happy   & 51,22\% & 41,06\% & \textcolor{red}{$-$10,16\%} \\
Neutral & 51,83\% & \textbf{56,81\%} & \textcolor{blue}{+4,98\%} \\
Sad     & 53,02\% & \textbf{58,14\%} & \textcolor{blue}{\textbf{+5,12\%}} \\
Angry   & 77,55\% & 66,67\% & \textcolor{red}{$-$10,88\%} \\
\midrule
\textbf{Overall} & \textbf{59,09\%} & \textbf{56,16\%} & \textcolor{red}{$-$2,93\%} \\
\bottomrule
\end{tabular}
\end{table}

Focal Loss với $\alpha_{Sad}{=}2{,}5$ chuyển năng lực học từ lớp dễ (Angry)
sang lớp khó (Sad, Neutral).
Trong ứng dụng thực tế (call center, sàng lọc tâm lý), phát hiện cảm xúc
tiêu cực quan trọng hơn tối đa hóa overall accuracy.

Hình~\ref{fig:confusion_final} trình bày confusion matrix của mô hình Final.

\begin{figure}[H]
  \centering
  \includegraphics[width=0.60\textwidth]{final_confusion_matrix.png}
  \caption{Normalized Confusion Matrix của mô hình Final
  (TME + BiGRU 256d + Focal Loss).}
  \label{fig:confusion_final}
\end{figure}

\subsection{Tổng hợp kết quả toàn bộ đề tài}

Bảng~\ref{tab:summary} tổng hợp kết quả qua các giai đoạn, cho thấy lộ
trình phát triển từ baseline đơn giản đến kiến trúc cuối cùng.

\begin{table}[H]
\centering
\caption{Tổng hợp kết quả qua các giai đoạn thực nghiệm.}
\label{tab:summary}
\begin{tabular}{@{}llrc@{}}
\toprule
\textbf{Giai đoạn} & \textbf{Mô hình} & \textbf{Test Acc.} & \textbf{Ghi chú} \\
\midrule
Midterm & A1: WavLM $\to$ MLP (audio-only)    & 56,44\% & Baseline \\
Midterm & B3b: Weighted fusion $\alpha{=}0{,}9$ & 57,95\% & Tốt nhất simple fusion \\
Final   & TME + BiGRU 768d + CE               & 59,09\% & Chưa có Focal Loss \\
Final   & TME + BiGRU 256d + Focal Loss       & 56,16\% & Cân bằng per-class \\
\bottomrule
\end{tabular}
\end{table}

% ===========================================================================
\section{Thảo luận}
\label{sec:discussion}

\paragraph{Text không hiệu quả với simple fusion, nhưng có giá trị qua TME.}
B2 (concat) chỉ cải thiện +0,13\% so với A1. Cảm xúc tiếng Việt nằm chủ
yếu ở prosody, không ở nội dung lời nói. Module TME giải quyết vấn đề này
bằng cosine gate (Công thức~\ref{eq:tme}), cho phép mô hình chọn lọc frame
text phù hợp thay vì trộn đều.

\paragraph{Model Compression là cần thiết cho low-resource.}
Giảm chiều từ 768 xuống 256 không gây mất mát đáng kể về overall accuracy
(59,09\% vs 56,16\%, trong đó phần lớn chênh lệch đến từ Focal Loss chứ
không phải compression). Trên tập 5.280 mẫu, mô hình nhỏ hơn generalize
tốt hơn nhờ tránh overfitting.

\paragraph{Focal Loss: đánh đổi có chủ đích.}
Giảm 2,93\% overall để tăng 5,12\% Sad là đánh đổi hợp lý trong ngữ cảnh
ứng dụng thực tế. Giá trị $\alpha$ có thể tiếp tục tinh chỉnh.

\paragraph{Giới hạn.}
Kết quả Final báo cáo trên 1 fold (fold 0, seed 42). Bản hoàn thiện cần
5-fold CV với mean $\pm$ std. Mô hình chưa được kiểm chứng trên tập ngoài
(ví dụ VNEMOS~\cite{vnemos2024}) do hạn chế thời gian thu thập dữ liệu.

% ===========================================================================
\section{Ứng dụng Web Demo}
\label{sec:demo}

Pipeline được đóng gói thành ứng dụng Web sử dụng Gradio~\cite{gradio2019}.
Người dùng ghi âm qua microphone hoặc upload file audio; hệ thống hiển thị
transcript (PhoWhisper) và phân phối xác suất cảm xúc (ViMamba-SER).
Nhờ compression xuống 256 chiều, mô hình chạy mượt trên CPU mà không cần GPU.

% ===========================================================================
\section{Kết luận}
\label{sec:conclusion}

Đồ án ViMamba-SER hoàn thành mục tiêu xây dựng hệ thống SER tiếng Việt
end-to-end, với các đóng góp chính:
(1)~Xây dựng baseline audio-only (56,44\%) và chứng minh text-only từ ASR
không hiệu quả cho SER tiếng Việt (44,32\%);
(2)~Chứng minh simple fusion không đủ (+1,51\%), cung cấp bằng chứng cho
sự cần thiết của module TME;
(3)~Thiết kế kiến trúc TME + BiGRU Fallback với Model Compression
(768$\to$256), giảm 80MB dung lượng và chống overfit;
(4)~Áp dụng Focal Loss tăng Sad +5,12\% và Neutral +4,98\%;
(5)~Đóng gói Web Demo real-time chứng minh tính khả thi triển khai.

\paragraph{Hướng phát triển.}
(1)~5-fold CV báo cáo mean $\pm$ std;
(2)~Cài đặt Bi-Mamba trên GPU có CUDA $\geq$ 11.8 để so sánh với BiGRU;
(3)~Mở rộng sang tập VNEMOS với 5 nhãn;
(4)~Thử PhoWhisper-Large để cải thiện transcript.

% ===========================================================================
\FloatBarrier
\clearpage
\bibliographystyle{splncs04}
\bibliography{references}

\end{document}
